% Generated by code/interviews-1/03-speed-drills/python/gen_drill.py. Do not edit. French copy by draft/pilot-quant-fr/ivgen_fr.py.
\ivfamily{Puissances de deux}{10, 23, 29, 30, 38, 45}{$2^{10} = 1024$. À partir de là, doubler ou diviser par deux, ou lire $2^{13}$ comme $8 \times 1024$.}
\ivfamily{Autres puissances}{6, 12, 17, 19, 21, 28, 34, 39, 44}{Élever au carré, puis multiplier une fois de plus si besoin~: $7^4 = 49^2$. $11^3 = 1331$ et $5^n$ vaut $10^n$ divisé par $2^n$.}
\ivfamily{Carrés jusqu'à 29}{4, 8, 9, 11, 22, 26, 36}{Utiliser $n^2 = (n - d)(n + d) + d^2$, ou simplement les connaître.}
\ivfamily{Racines carrées}{2, 3, 13, 16, 24, 35, 41, 42}{La taille donne le chiffre des dizaines, et le chiffre des unités laisse au plus deux choix. Comparer avec le carré du nombre finissant par 5.}
\ivfamily{Racines cubiques}{14, 25, 31, 32, 43}{Connaître les cubes jusqu'à 15. Le chiffre des unités d'un cube fixe celui de sa racine.}
\ivfamily{Logarithmes}{5, 7, 15, 27, 37}{Réécrire le nombre comme une puissance de la base~: $\log_2 4096$ est la puissance de deux qui donne 4096.}
\ivfamily{Factorielles}{1, 18, 20, 33, 40}{Connaître jusqu'à $9! = 362\,880$. Un rapport tel que $\tfrac{10!}{8!}$ vaut simplement $10 \times 9$.}
